% SEGMENT OLP-0053-S01
% Part: sets-functions-relations
% Chapter: infinite
% Section: dedekinds-proof
%
\documentclass[../../../include/open-logic-section]{subfiles}

\begin{document}

\olfileid{sfr}{infinite}{dedekindsproof}

\olsection[டெடெகிண்டின் ``நிறுவல்'']{முடிவுறாத கணமொன்று இருப்பதற்கான
டெடெகிண்டின் ``நிறுவல்''}

இந்த அத்தியாயத்தில், டெடெகிண்ட் இயற்கணிதங்களின் வழியாக இயல் எண்களைக்
கணக்கோட்பாட்டு முறையில் கையாண்டோம். \olref[arith][ref]{sec} இல்,
பகுப்பாய்வியலின் எண்முறைமைக் கட்டமைப்பின் மெய்யியல் முக்கியத்துவத்தைப்
(பிறவற்றுடன்) பற்றி சிந்தித்தோம். இங்கு நாம் சாதித்திருப்பதன்
முக்கியத்துவத்தை இப்போது சிந்திக்க வேண்டும்.

% SEGMENT OLP-0053-S02
\olref[sfr][arith][]{chap} முழுவதிலும், இயல் எண்களை நமக்குத்
தரப்பட்டவையாக ஏற்றுக்கொண்டு, அவற்றைப் பயன்படுத்தி முழுக்கள்,
விகிதமுறு எண்கள், மெய்யெண்கள் ஆகியவற்றை வெளிப்படையாகக் கட்டமைத்தோம்.
இந்த அத்தியாயத்தில், இயல் எண்களை வெளிப்படையாகக் கட்டமைக்கவில்லை.
\emph{எந்த டெடெகிண்ட்-முடிவுறாத கணமும் தரப்பட்டால்}, $\Nat$ எவ்வாறு
செயல்பட வேண்டும் என்று நாம் விரும்புகிறோமோ அவ்வாறே செயல்படும் ஒரு
கணத்தை வரையறுக்கலாம் என்று மட்டுமே காட்டியுள்ளோம்.

% SEGMENT OLP-0053-S03
எனவே, \emph{எந்தப் பொருள்களே இயல் எண்கள்} என்பது போன்ற ஒரு
மெய்ப்பொருளியல் வினாவுக்கு விடையளித்துவிட்டோம் என்று நாம் கூற முடியாது
என்பது தெளிவு. அது நல்லதும்கூட. ஏனெனில், $\Real$ ஐ டெடெகிண்ட்
வெட்டுகளின் கணம் என்று வரையறுப்பதை ஒரு வசதியான விதிப்பாகக் கொள்ளாமல்,
ஒரு \emph{கண்டுபிடிப்பாகக்} கொள்வது தவறு என்று
\olref[sfr][arith][ref]{sec} இல் வலியுறுத்தினோம்.
மெய்யெண்களின் முதன்மையான கணிதப் பண்புகளுக்கு டெடெகிண்ட் வெட்டுகள்
எடுத்துக்காட்டுகளாக அமைகின்றன என்பதே தீர்மானமான கவனிப்பு.
இங்கும் அவ்வாறே: இயல் எண்களின் முதன்மையான கணிதப் பண்புகளுக்கு
\emph{எந்த} டெடெகிண்ட் இயற்கணிதமும் எடுத்துக்காட்டாக அமைகிறது என்பதே
தீர்மானமான கவனிப்பு. (உண்மையில், எல்லா டெடெகிண்ட் இயற்கணிதங்களும்
\emph{சமவடிவமானவை} என்று நிறுவி, டெடெகிண்ட் இந்தக் கருத்தை
வலுவுறுத்தினார் (\citeyear[Theorems
132--3]{Dedekind1888}).
எனவே, தற்கால ``அமைப்பியலாளர்கள்'' பலர் டெடெகிண்டை ஒரு முன்னோடியாகக்
குறிப்பிடுவதில் வியப்பில்லை.)

% SEGMENT OLP-0053-S04
மேலும், இன்னும் ஓரளவு நெறிப்படுத்தப்படாததாகவே உள்ள, ஆனால் இயல்
எண்களை நமக்குத் தரப்பட்டவையாகக் (வெளிப்படையாக) கருதாத முறைசாராத
எளிய கணக்கோட்பாட்டிற்குள் இயல் எண்களின் கோட்பாட்டை எவ்வாறு பதிக்கலாம்
என்பதைக் காட்டியுள்ளோம். ஆகவே, டெடெகிண்ட் பின்வருமாறு விளக்கிய அவரது
சொந்த உயர்நோக்குத் திட்டத்தை நிறைவேற்றும் பாதையில் நாம் இருக்கலாம்:

% SEGMENT OLP-0053-S05
\begin{quote}
	அறிவியலில் நிறுவக்கூடிய எதையும் நிறுவல் இல்லாமல் நம்பக் கூடாது.
	இந்தக் கோரிக்கை நியாயமானதாகத் தோன்றினாலும், மிக எளிய அறிவியலுக்கு
	அடித்தளம் அமைக்கும் அண்மைக்கால முறைகளில்கூட அது நிறைவேற்றப்பட்டுள்ளது
	என்று என்னால் கருத முடியாது; அதாவது, எண்களின் கோட்பாட்டைக் கையாளும்
	தருக்கவியலின் அந்தப் பகுதியில். எண்கணிதத்தை (இயற்கணிதம்,
	பகுப்பாய்வியல்) தருக்கவியலின் ஒரு பகுதி மட்டுமே என்று கூறுவதன் மூலம்,
	எண்-கருத்தை வெளி, காலம் பற்றிய கருத்துகளிலிருந்தோ உள்ளுணர்வுகளிலிருந்தோ
	முற்றிலும் சார்பற்றதாக நான் கருதுகிறேன் என்பதையே உணர்த்துகிறேன்---
	அதனைத் தூய சிந்தனை விதிகளின் நேரடி விளைபொருளாகவே கருதுகிறேன்.
	\citep[preface]{Dedekind1888}
\end{quote}

% SEGMENT OLP-0053-S06
டெடெகிண்டின் துணிவான கருத்து இதுதான். முறைசாராத கணக்கோட்பாட்டை
மட்டுமே பயன்படுத்தி இயல் எண்களை எவ்வாறு அமைப்பது என்று இப்போதுதான்
காட்டினோம். இயல் எண்களும் சிறிது கணக்கோட்பாடும் தரப்பட்டால்
மெய்யெண்களை எவ்வாறு கட்டமைப்பது என்று \olref[sfr][arith][]{chap} இல்
கண்டோம். ஆகவே, ஒருவேளை, ``எண்கணிதம் (இயற்கணிதம், பகுப்பாய்வியல்)''
என்பது ``தருக்கவியலின் ஒரு பகுதி மட்டுமே'' என்று திகழலாம்
(டெடெகிண்ட் ``தருக்கவியல்'' என்ற சொல்லுக்கு அளித்த விரிவான பொருளில்).

% SEGMENT OLP-0053-S07
கருத்து இதுதான். ஆனால் சற்று நிறுத்துங்கள். ஒரு டெடெகிண்ட் இயற்கணிதத்தை
(இயல் எண்களுக்கான நமது பதிலியை) நாம் கட்டமைப்பது, ஒரு
டெடெகிண்ட்-முடிவுறாத கணம் இருப்பதை நிபந்தனையாகக் கொண்டுள்ளது.
(\olref[sfr][infinite][dedekind]{thm:DedekindInfiniteAlgebra} ஐத்
திரும்பிப் பாருங்கள்.) ஒரு டெடெகிண்ட்-முடிவுறாத கணம் இருப்பதைத்
``தருக்கவியல்'' அல்லது ``தூய சிந்தனை விதிகள்'' கொண்டு நிறுவ
முடியாவிட்டால், இந்தத் திட்டம் நின்றுவிடும்.

% SEGMENT OLP-0053-S08
அப்படியானால், ஒரு டெடெகிண்ட்-முடிவுறாத கணம் இருப்பதைத்
``தூய சிந்தனை விதிகள்'' கொண்டு நிறுவ \emph{முடியுமா}?
டெடெகிண்டின் முயற்சி இதோ:

% SEGMENT OLP-0053-S09
\begin{quote}
  என் சொந்தச் சிந்தனைகளின் மண்டலம், அதாவது என் சிந்தனைக்குப்
  பொருள்களாக இருக்கக்கூடிய அனைத்துப் பொருள்களின் ஒட்டுமொத்தம் $S$,
  முடிவுறாதது. ஏனெனில் $s$ என்பது $S$ இன் ஓர் உறுப்பு எனில்,
  $s$ என் சிந்தனைக்குப் பொருளாக இருக்க முடியும் என்ற சிந்தனை $s'$
  ஆனதும் $S$ இன் ஓர் உறுப்பாகும். இதனை உறுப்பு $s$ இன் படிமம்
  $\phi(s)$ எனக் கொண்டால், அப்போது \dots~$S$ ஆனது [டெடெகிண்ட்]
  முடிவுறாதது; நிறுவ வேண்டியது இதுவே.
	\citep[\S66]{Dedekind1888}
\end{quote}

% SEGMENT OLP-0053-S10
மிகவும் கடுமையான கணித நிறுவல்களையே பெரும்பாலும் கொண்ட ஒரு நூலின்
நடுவில் இதனைக் காண்பது மிகவும் வியப்பூட்டுகிறது. இரண்டு குறிப்புகளைக்
கூறுவது பயனுள்ளது.

% SEGMENT OLP-0053-S11
முதலாவது: இந்த ``நிறுவல்'' இன்று நாம் ``கணிதத்தன்மை'' என்று
அடையாளம் காணக்கூடிய பண்பை ஏறத்தாழக் கொண்டிருக்கவில்லை.
இது உளவியல் பொருள்களைப் (சிந்தனைகளைப்) பேசுகிறது; அதிலும் வெறும்
\emph{இருக்கக்கூடிய} பொருள்களையே பேசுகிறது.

% SEGMENT OLP-0053-S12
இரண்டாவது: டெடெகிண்ட் இயற்கணிதங்களை நாம் வழங்கியுள்ள விதத்திலாவது,
இந்த ``நிறுவலில்'' ஒரு நேரடியான தொழில்நுட்பக் குறைபாடு உள்ளது.
டெடெகிண்டின் வாதம் வெற்றிபெற்றால், முடிவுறாத அளவு பல பொருள்கள்
(குறிப்பாக, முடிவுறாத அளவு பல சிந்தனைகள்) உள்ளன என்று மட்டுமே
அது நிறுவுகிறது. ஆனால் $S$ ஐ, பன்மையில் $S$ என்பதை
\emph{சில பொருள்கள்} என்று
கருதாமல், முடிவுறாத அளவு பல !!{element}s கொண்ட ஒரே \emph{கணமாகக்}
கருதுவதற்கான காரணத்தையும் டெடெகிண்ட் நமக்குத் தர வேண்டும்.

% SEGMENT OLP-0053-S13
டெடெகிண்ட் இங்கு ஒரு இடைவெளியைக் காணவில்லை என்ற உண்மை, அவர்
``ஒட்டுமொத்தம்'' என்ற சொல்லைப் பயன்படுத்திய விதம், \emph{நாம்} ``கணம்'' என்ற
சொல்லைப் பயன்படுத்தும் விதத்தைத் துல்லியமாகப் பின்பற்றவில்லை என்று
உணர்த்தலாம்.\footnote{அவ்வாறு பின்பற்றவில்லை என்று எண்ணுவதற்கு வேறு
காரணங்களும் நமக்கு உள்ளன; \citet[p.~23]{Potter2004} ஐப் பார்க்கவும்.}
ஆனால் இது மிகுந்த வியப்பிற்குரியதன்று. கடந்த இரண்டு அத்தியாயங்களில்
நாம் மேற்கொண்ட திட்டம்---இயல் எண்களின் ஒரு ``கட்டமைப்பு'', அவற்றிலிருந்து
முழுக்கள், மெய்யெண்கள், விகிதமுறு எண்கள் ஆகியவற்றின் ஒரு
``கட்டமைப்பு''---முழுவதும் முறைசாராத விதத்தில் மேற்கொள்ளப்பட்டது.
எந்தக் கணங்கள் எப்போது உள்ளன என்பது பற்றிய பல பொதுக் கொள்கைகளை
வகுக்காமலேயே, நமக்குத் தேவைப்பட்டபோதெல்லாம் இந்தக் கணம் அல்லது அந்தக்
கணம் உள்ளதாக எடுத்துக்கொண்டோம். ஆனால் \emph{சில} பொதுக் கொள்கைகள்
நமக்குத் தேவை என்பது தெரியும்; இல்லாவிட்டால் ரசலின் முரண்பாட்டில்
விழுந்துவிடுவோம்.

நமது முறைசாராமையைக் கடந்து வளர வேண்டிய நேரம் வந்துவிட்டது.

\end{document}
